% 【本文件写什么】impl 实现层：task
% - 定位：对 wateros_task::WaitQueue 的薄包装。
% - crate 路径：os/components/wateros-ipc/ipc-waitqueue/waitqueue-impl/impl-task/src/
% - 如何实现对应 api-v0 trait：关键类型、状态机、数据结构
% - 算法或硬件交互流程，可用伪代码/流程图
% - 本 impl 启用的 Cargo [features] 与 arch feature，impl-riscv64 等
% - 与其它 impl 变体的差异、切换 feature 时的影响
% - 性能/内存假设与已知限制
% 【事实来源】
% - 上述 crate 源码与 Cargo.toml
% - os/feature-tree.txt 中指向本 impl 的 feature 链

\subsubsection{impl-task}

\code{impl-task}将\code{wateros\_task::wait\_queue::WaitQueue}包装为\code{WaitQueue}，字段仅含\code{inner}，不在 IPC 层缓存 waiters 或附加锁。

\begin{rustcode}
#[derive(Clone, Copy, Debug, PartialEq, Eq)]
pub struct WaitQueue {
    inner: TaskWaitQueue,
}
\end{rustcode}

全部\code{wait\_*}、\code{wake\_*}方法内联委托\code{inner}；\code{try\_release\_empty}在队列无等待者时释放底层\code{WaitQueueId}供复用，futex 空键清理依赖此路径。\code{requeue\_to}将部分等待者迁移到目标队列，支撑\code{FUTEX\_REQUEUE}语义，未列入\code{IpcWaitQueueOps} trait 以避免 v0 契约膨胀。

\code{WaitQueue}实现\code{IpcWaitQueueOps}，同时作为\code{Copy}句柄在\code{FutexHub}的\code{BTreeMap<FutexKey, WaitQueue>}中存储。性能假设与 task 层相同：单核 bring-up 下由调度器保证唤醒顺序，无额外 IPC 层开销。
